% !TeX root = ../main.tex
\section{Finite automata}
\label{sec:automata}

A grammar generates the words of a language. An automaton goes the other way:
it reads a word and decides whether it belongs to the language. This section
follows Hopcroft, Motwani and Ullman~\cite[Ch.~2]{hopcroft2007automata}. The
lecture writes automata in the Gold language; here they are drawn as diagrams
and, when they need to run, written in Rascal.

\subsection{Labeled transition systems}

\begin{definition}
A \textbf{labeled transition system} (LTS) is a triple $\mathcal{M} = \langle Q,
\Sigma, \delta \rangle$ where $Q$ is the set of states of the system, $\Sigma$ is
the set of actions (labels) and $\delta \subseteq Q \times \Sigma \times Q$ is the
set of transitions. A transition $\langle q, \sigma, q' \rangle$ says that the
system can move from $q$ to $q'$ when the action $\sigma$ happens.
\end{definition}

An automaton is an LTS with a starting point and a notion of success. The
lecture motivates it with a purchase. A customer pays the store; the store
ships the goods; the store redeems the money at the bank; the bank transfers
the money to the store. Shipping can happen before or after the money is
redeemed or transferred, but the purchase is complete only when the goods
have been shipped and the money transferred. Figure~\ref{fig:shopping} models
this.

\begin{figure}[h]
\centering
\begin{tikzpicture}[node distance=2.3cm, auto]
  \node[state, initial above] (A) {A};
  \node[state, below=1.4cm of A] (B) {B};
  \node[state, right=of B] (D) {D};
  \node[state, right=of D] (F) {F};
  \node[state, below=1.6cm of B] (C) {C};
  \node[state, right=of C] (E) {E};
  \node[state, accepting, right=of E] (G) {G};
  \path[->]
    (A) edge node {pay} (B)
    (B) edge node {redeem} (D)
    (D) edge node {transfer} (F)
    (B) edge node[swap] {ship} (C)
    (D) edge node {ship} (E)
    (F) edge node {ship} (G)
    (C) edge node[swap] {redeem} (E)
    (E) edge node[swap] {transfer} (G);
\end{tikzpicture}
\caption{The purchase automaton. $\Sigma = \{\mathit{pay}, \mathit{ship},
\mathit{redeem}, \mathit{transfer}\}$, the initial state is A and the only
final state is G.}
\label{fig:shopping}
\end{figure}

Every path from A to G is a correct execution of the purchase. The word
$\mathit{pay}\ \mathit{ship}\ \mathit{redeem}\ \mathit{transfer}$ reaches G, and so
does $\mathit{pay}\ \mathit{redeem}\ \mathit{transfer}\ \mathit{ship}$. The word
$\mathit{pay}\ \mathit{ship}\ \mathit{ship}$ does not: the goods cannot be
shipped twice, so C has no transition for $\mathit{ship}$. In this example the
symbols of the alphabet are whole actions, not characters.

\subsection{Deterministic finite automata}

\begin{definition}
A \textbf{deterministic finite automaton} (DFA) is a 5-tuple $\mathcal{A} =
\langle Q, \Sigma, q_I, F, \delta \rangle$ where $Q$ is a finite set of states,
$\Sigma$ is the input alphabet, $q_I \in Q$ is the initial state, $F \subseteq Q$
is the set of final states and $\delta : Q \times \Sigma \to Q$ is the transition
function.
\end{definition}

The automaton reads the word from an input tape, one symbol at a time, and
each symbol moves it to the next state.

\begin{definition}
A word $\omega = \sigma_1 \cdots \sigma_n \in \Sigma^*$ is \textbf{accepted} by
$\mathcal{A}$ if there are states $q_0, \dots, q_n$ such that $q_0 = q_I$,
$q_n \in F$ and $\delta(q_{i-1}, \sigma_i) = q_i$ for every $i \in \{1, \dots, n\}$.
\end{definition}

The same idea can be packed into one function that processes a whole word.
The \textbf{extended transition function} $\hat\delta : Q \times \Sigma^* \to Q$
is defined by induction on the word:
\[
  \hat\delta(q, \lambda) = q
  \qquad\qquad
  \hat\delta(q, \omega\sigma) = \delta(\hat\delta(q, \omega), \sigma).
\]
The empty word leaves the automaton where it is, and a word that ends in
$\sigma$ is processed by first processing the prefix and then taking one more
step.

\begin{definition}
The \textbf{language of a DFA} is $L(\mathcal{A}) = \{ \omega \in \Sigma^* \mid
\hat\delta(q_I, \omega) \in F \}$.
\end{definition}

\paragraph{Example: words that contain 01.} The lecture builds a DFA over
$\{0, 1\}$ that accepts the words containing the substring $01$
(Figure~\ref{fig:contains01}). The state names say what the automaton has
seen. In \textit{Just1} it has seen only 1s (or nothing). A 0 moves it to
\textit{Read0}, where further 0s change nothing because the last symbol is
still 0. A 1 after that completes the pattern, and from \textit{Read01} every
symbol keeps the word accepted.

\begin{figure}[h]
\centering
\begin{tikzpicture}[node distance=3cm, auto, every state/.style={draw, minimum size=11mm, font=\small}]
  \node[state, initial] (J) {\textit{Just1}};
  \node[state, right=of J] (R0) {\textit{Read0}};
  \node[state, accepting, right=of R0] (R01) {\textit{Read01}};
  \path[->]
    (J) edge[loop above] node {1} ()
    (J) edge node {0} (R0)
    (R0) edge[loop above] node {0} ()
    (R0) edge node {1} (R01)
    (R01) edge[loop above] node {0, 1} ();
\end{tikzpicture}
\caption{DFA for the words over $\{0,1\}$ that contain $01$.}
\label{fig:contains01}
\end{figure}

For the word $1010$:
$\hat\delta(\mathit{Just1}, 1010) = \delta(\delta(\delta(\delta(\mathit{Just1}, 1), 0), 1), 0)
= \delta(\delta(\delta(\mathit{Just1}, 0), 1), 0) = \delta(\delta(\mathit{Read0}, 1), 0)
= \delta(\mathit{Read01}, 0) = \mathit{Read01}$, so $1010$ is accepted.
Table~\ref{tab:contains01} gives the other words of the lecture.

\begin{table}[h]
\centering
\small
\begin{tabular}{l l c}
  \toprule
  Word & States visited & Accepted \\
  \midrule
  $01$     & Just1, Read0, Read01                          & yes \\
  $1010$   & Just1, Just1, Read0, Read01, Read01           & yes \\
  $100011$ & Just1, Just1, Read0, Read0, Read0, Read01, Read01 & yes \\
  $0$      & Just1, Read0                                  & no \\
  $111000$ & Just1, Just1, Just1, Just1, Read0, Read0, Read0 & no \\
  \bottomrule
\end{tabular}
\caption{Runs of the automaton of Figure~\ref{fig:contains01}.}
\label{tab:contains01}
\end{table}

\paragraph{Automata in Rascal.} In the code for this section an automaton is a
set of initial states, a set of final states and a relation of transitions.
The empty string \texttt{""} is used as the label of a $\lambda$-transition,
which is needed in the next subsection. A DFA is the case with one initial
state and one transition per state and symbol.
\coderange{Rascal}{rascal/src/Automata.rsc}{types}
The automaton of Figure~\ref{fig:contains01} is written as
\coderange{Rascal}{rascal/src/Automata.rsc}{contains01}
and \texttt{Automata.rsc} checks it against the words of
Table~\ref{tab:contains01} and against every word of length up to 8.

\subsection{Nondeterministic finite automata}

DFA definitions are restrictive: every state needs exactly one transition for
every symbol, and there is one initial state. A nondeterministic automaton
relaxes three things.
\begin{enumerate}
  \item Transitions may be over words, including the empty word $\lambda$.
    A transition labeled $abb$ reads three symbols; a transition labeled
    $\lambda$ reads nothing.
  \item There is choice: a state may have several transitions with the same
    label, or none.
  \item There may be several initial states.
\end{enumerate}

\begin{definition}
A \textbf{nondeterministic finite automaton} (NFA) is a 5-tuple $\mathcal{A} =
\langle Q, \Sigma, I, F, \Delta \rangle$ where $I \subseteq Q$ is a set of
initial states and $\Delta \subseteq Q \times (\Sigma \cup \{\lambda\}) \times Q$
is a transition relation. Equivalently, $\Delta$ is a function from $Q \times
(\Sigma \cup \{\lambda\})$ to $\mathcal{P}(Q)$, the set of subsets of $Q$.
\end{definition}

\begin{definition}
A word $\omega = \sigma_1 \cdots \sigma_n$ is \textbf{recognized} by an NFA if
there is a path from some initial state $q_I \in I$ to some final state
$f \in F$ that reads all the symbols of $\omega$, in order, possibly taking
$\lambda$-transitions in between. We write $q_I \stackrel{\omega}{\Longrightarrow} f$.
\end{definition}

The word is accepted if one path works; other paths may get stuck or end in
non-final states. A DFA is an NFA with $I = \{q_I\}$, no $\lambda$-transitions
and exactly one transition per state and symbol.

\paragraph{Example: last symbol 0, or only a 1.} The lecture asks for an
automaton over $\{0,1\}$ that recognizes the words whose last symbol is 0 and
the word $1$. With $\lambda$-transitions the two cases can be built
separately and joined at the start (Figure~\ref{fig:last0}). From A the
automaton chooses, without reading anything, which case it is in. The
branch through B reads exactly one 1. The branch through C reads any word and
then guesses that the current 0 is the last symbol. If the guess is wrong,
that path gets stuck in D, which has no outgoing transitions.

\begin{figure}[h]
\centering
\begin{tikzpicture}[node distance=2.2cm, auto]
  \node[state, initial] (A) {A};
  \node[state, above right=0.8cm and 2cm of A] (B) {B};
  \node[state, below right=0.8cm and 2cm of A] (C) {C};
  \node[state, accepting, below right=0.8cm and 2cm of B] (D) {D};
  \path[->]
    (A) edge node {$\lambda$} (B)
    (A) edge node[swap] {$\lambda$} (C)
    (B) edge node {1} (D)
    (C) edge node[swap] {0} (D)
    (C) edge[loop below] node {0, 1} ();
\end{tikzpicture}
\caption{NFA for the words that end in 0, plus the word 1.}
\label{fig:last0}
\end{figure}

\paragraph{Example: \texttt{cat} and \texttt{hat}.} Over $\{c, a, t, h\}$ the
lecture defines the NFA of Figure~\ref{fig:cathat}. It has two branches
joined by $\lambda$-transitions from A. The upper branch loops in B on every
symbol, guesses where \texttt{hat} starts, reads it and stays in E. It
accepts the words that contain \texttt{hat}. The lower branch is a DFA that
remembers how much of \texttt{cat} it has just read: G after \texttt{c}, H
after \texttt{ca}. A \texttt{t} in H means the word contains \texttt{cat}, and
I is a trap state that is not final. F, G and H are final, so the lower branch
accepts the words that do not contain \texttt{cat}.

\begin{figure}[h]
\centering
\begin{tikzpicture}[node distance=2.1cm, auto, every state/.style={draw, minimum size=7mm, inner sep=1pt, font=\small}]
  \node[state, initial] (A) {A};
  \node[state, above right=1.2cm and 1.6cm of A] (B) {B};
  \node[state, right=of B] (C) {C};
  \node[state, right=of C] (D) {D};
  \node[state, accepting, right=of D] (E) {E};
  \node[state, accepting, below right=1.2cm and 1.6cm of A] (F) {F};
  \node[state, accepting, right=of F] (G) {G};
  \node[state, accepting, right=of G] (H) {H};
  \node[state, right=of H] (I) {I};
  \path[->]
    (A) edge node {$\lambda$} (B)
    (A) edge node[swap] {$\lambda$} (F)
    (B) edge[loop above] node {c, a, t, h} ()
    (B) edge node {h} (C)
    (C) edge node {a} (D)
    (D) edge node {t} (E)
    (E) edge[loop above] node {c, a, t, h} ()
    (F) edge[loop below] node {a, t, h} ()
    (F) edge[bend left=15] node {c} (G)
    (G) edge[bend left=15] node {t, h} (F)
    (G) edge[loop above] node {c} ()
    (G) edge[bend left=15] node {a} (H)
    (H) edge[bend left=15] node {c} (G)
    (H) edge[bend left=50] node {a, h} (F)
    (H) edge node {t} (I)
    (I) edge[loop below] node {c, a, t, h} ();
\end{tikzpicture}
\caption{The NFA of the lecture over $\{c,a,t,h\}$.}
\label{fig:cathat}
\end{figure}

Together, the automaton accepts the words that do not contain \texttt{cat}
or that contain \texttt{hat}. The slide states the exercise as ``does not
contain the string cat or the string hat'', which can also be read as
``contains neither''. Under that reading the automaton is wrong, since it
accepts \texttt{hat} and \texttt{cathat}. Exercise~\ref{ex:neither} builds
the automaton for the second reading. In Rascal:
\coderange{Rascal}{rascal/src/Automata.rsc}{cathat}
Running an NFA means keeping the set of states where some path could be.
Before reading and after every symbol, the set is closed under
$\lambda$-transitions:
\coderange{Rascal}{rascal/src/Automata.rsc}{run}

\subsection{From NFA to DFA}

Nondeterminism makes automata easier to write, but it does not let them
recognize more languages.

\begin{theorem}
Every NFA has an equivalent DFA.
\end{theorem}

The construction removes the three relaxations one at a time. All three use
the same idea: a state of the new automaton stands for a set of states of the
old one, and it behaves like all of them together. A set is initial or final
according to the states it contains.

\paragraph{Transitions over words.} A transition $\langle q, abb, q' \rangle$
becomes a chain $q \xrightarrow{a} p_1 \xrightarrow{b} p_2 \xrightarrow{b} q'$
through new states $p_1, p_2$.

\paragraph{Eliminate $\lambda$-transitions.} The \textbf{$\lambda$-closure}
$\mathit{cl}_\lambda(q)$ is the set of states reachable from $q$ using only
$\lambda$-transitions, including $q$ itself. If the automaton can go from $u$
to $q$ reading $a$, and $t \in \mathit{cl}_\lambda(q)$, then it can also go from
$u$ to $t$ reading $a$: add $\langle u, a, t \rangle$. A state becomes final if
its closure contains a final state, and the initial states are replaced by
their closures. After that, the $\lambda$-transitions can be deleted.

\paragraph{Eliminate choice.} If a state has several transitions with the same
symbol, replace the targets by one state that stands for the set of them.
This is the \textbf{set construction}. Starting from the set of initial
states, for each set $S$ and symbol $\sigma$ the next set is
\[
  \delta'(S, \sigma) = \bigcup_{q \in S} \mathit{cl}_\lambda\big(\{ t \mid \langle q, \sigma, t \rangle \in \Delta \}\big),
\]
and $S$ is final if $S \cap F \ne \emptyset$. Only the sets reachable from the
start are generated. The empty set appears when no state has a transition for
$\sigma$; it is a trap state.

\paragraph{Remove additional initial states.} The set $I$ of initial states
becomes a single initial state of the new automaton. In the set construction
this happens by starting from $\mathit{cl}_\lambda(I)$ instead of from one
state.

In Rascal the whole construction is a search over sets of states. Each set
is named by listing its states, so \texttt{\{A,S\}} is a state of the DFA:
\coderange{Rascal}{rascal/src/Automata.rsc}{determinize}

\paragraph{Example.} Exercise~6 of the lecture asks to convert the NFA on the
left of Figure~\ref{fig:ex6} to a DFA. S reads $a$ and may stay in S or go to
A; it reads $b$ and may go to B or C. Table~\ref{tab:ex6} runs the set
construction from $\{S\}$, and the DFA is on the right of the figure.

\begin{table}[h]
\centering
\small
\begin{tabular}{l l l c}
  \toprule
  Set & on $a$ & on $b$ & Final \\
  \midrule
  $\{S\}$       & $\{A, S\}$    & $\{B, C\}$ & no \\
  $\{A, S\}$    & $\{A, C, S\}$ & $\{B, C\}$ & no \\
  $\{B, C\}$    & $\{C\}$       & $\emptyset$ & yes \\
  $\{A, C, S\}$ & $\{A, C, S\}$ & $\{B, C\}$ & yes \\
  $\{C\}$       & $\{C\}$       & $\emptyset$ & yes \\
  $\emptyset$   & $\emptyset$   & $\emptyset$ & no \\
  \bottomrule
\end{tabular}
\caption{Set construction for the NFA of Figure~\ref{fig:ex6}.}
\label{tab:ex6}
\end{table}

\begin{figure}[h]
\centering
\begin{tikzpicture}[node distance=1.8cm, auto]
  \node[state, initial] (S) {S};
  \node[state, right=2.2cm of S] (A) {A};
  \node[state, accepting, below=1.6cm of A] (C) {C};
  \node[state, below=1.6cm of S] (B) {B};
  \path[->]
    (S) edge[loop above] node {a} ()
    (S) edge node {a} (A)
    (S) edge node {b} (C)
    (S) edge node[swap] {b} (B)
    (A) edge node {a} (C)
    (B) edge node[swap] {a} (C)
    (C) edge[loop right] node {a} ();
\end{tikzpicture}
\hspace{1cm}
\begin{tikzpicture}[node distance=2.2cm, auto, every state/.style={draw, minimum size=10mm, inner sep=1pt, font=\scriptsize}]
  \node[state, initial] (s) {$\{S\}$};
  \node[state, right=of s] (as) {$\{A,S\}$};
  \node[state, accepting, right=of as] (acs) {$\{A,C,S\}$};
  \node[state, accepting, below=1.5cm of as] (bc) {$\{B,C\}$};
  \node[state, accepting, right=of bc] (c) {$\{C\}$};
  \node[state, below=1.3cm of bc] (e) {$\emptyset$};
  \path[->]
    (s) edge node {a} (as)
    (as) edge node {a} (acs)
    (acs) edge[loop above] node {a} ()
    (s) edge node[swap] {b} (bc)
    (as) edge node[pos=0.4] {b} (bc)
    (acs) edge node[pos=0.3] {b} (bc)
    (bc) edge node {a} (c)
    (c) edge[loop right] node {a} ()
    (bc) edge node {b} (e)
    (c) edge node {b} (e)
    (e) edge[loop left] node {a, b} ();
\end{tikzpicture}
\caption{The NFA of exercise 6 (left) and the DFA produced by the set construction (right).}
\label{fig:ex6}
\end{figure}

Six sets appear, out of the $2^4 = 16$ subsets of $\{S, A, B, C\}$. In the
worst case an NFA with $n$ states needs a DFA with $2^n$ states, but usually
only a few sets are reachable. The tests in \texttt{Automata.rsc} check that
the construction produces exactly these six states and that both automata
accept the same words.

\subsection{Combining automata}

Automata for simple languages can be combined into automata for more complex
ones. The lecture uses two automata over $\{a, b\}$ as running examples
(Figure~\ref{fig:lb-l2a}). $L_b$ recognizes the words that contain a $b$.
$L_{2a}$ recognizes the words that end in $aa$: A loops on every symbol and
guesses where the final $aa$ starts. The $\lambda$ back from C lets the
automaton continue after an $aa$, but since only C is final, the word must
end there.

\begin{figure}[h]
\centering
\begin{tikzpicture}[node distance=2.6cm, auto]
  \node[state, initial] (W) {W};
  \node[state, accepting, right=of W] (Z) {Z};
  \node[above=0.9cm of W, xshift=-1.2cm] {$L_b$};
  \path[->]
    (W) edge[loop above] node {a, b} ()
    (W) edge node {b} (Z)
    (Z) edge[loop above] node {a, b} ();
\end{tikzpicture}
\hspace{1.2cm}
\begin{tikzpicture}[node distance=2.2cm, auto]
  \node[state, initial] (A) {A};
  \node[state, right=of A] (B) {B};
  \node[state, accepting, right=of B] (C) {C};
  \node[above=0.9cm of A, xshift=-1.2cm] {$L_{2a}$};
  \path[->]
    (A) edge[loop above] node {a, b} ()
    (A) edge[bend left=20] node {a} (B)
    (B) edge[bend left=20] node {b} (A)
    (B) edge node {a} (C)
    (C) edge[bend left=40] node {$\lambda$} (A);
\end{tikzpicture}
\caption{The two automata used for the combinations.}
\label{fig:lb-l2a}
\end{figure}

For $\mathcal{A}_1 = \langle Q_1, \Sigma, I_1, F_1, \Delta_1 \rangle$ and
$\mathcal{A}_2 = \langle Q_2, \Sigma, I_2, F_2, \Delta_2 \rangle$ with disjoint
sets of states:

\paragraph{Union.} Add a new initial state $q$ with a $\lambda$-transition to
every initial state of both automata:
\[
\mathcal{A}_1 \cup \mathcal{A}_2 = \langle Q_1 \cup Q_2 \cup \{q\}, \Sigma, \{q\}, F_1 \cup F_2, \Delta \rangle,
\quad
\Delta = \Delta_1 \cup \Delta_2 \cup \{ \langle q, \lambda, i \rangle \mid i \in I_1 \cup I_2 \}.
\]
From $q$ the automaton chooses which of the two to run.

\paragraph{Concatenation.} Connect every final state of the first automaton
to every initial state of the second:
\[
\mathcal{A}_1\mathcal{A}_2 = \langle Q_1 \cup Q_2, \Sigma, I_1, F_2, \Delta \rangle,
\quad
\Delta = \Delta_1 \cup \Delta_2 \cup \{ \langle f, \lambda, i \rangle \mid f \in F_1,\ i \in I_2 \}.
\]
For $L_b L_{2a}$ the automaton reads a prefix that contains a $b$, jumps from
Z to A and then reads a suffix that ends in $aa$.

\paragraph{Kleene star.} The star must accept the empty word and any sequence
of words of $L(\mathcal{A}_1)$. The slide gives
\[
\mathcal{A}_1^* = \langle Q_1, \Sigma, I_1, F_1 \cup I_1, \Delta \rangle,
\quad
\Delta = \Delta_1 \cup \{ \langle f, \lambda, i \rangle \mid f \in F_1,\ i \in I_1 \}
\]
(the slide writes $F_2$ and $\Delta_2$, but there is only one automaton). The
$\lambda$-transitions from the final states back to the initial states allow
repetition. Making the initial states final, to accept $\lambda$, does not work
in general. In $L_b$ the initial state W has a loop on $a$, so the word $a$
would end in W and be accepted, but $a \notin L_b^*$: it is not empty and it
contains no $b$. The safe construction adds a new initial state $q$ that is
final, with $\lambda$-transitions from $q$ to the old initial states and from the
old final states back to $q$:
\[
\mathcal{A}_1^* = \langle Q_1 \cup \{q\}, \Sigma, \{q\}, F_1 \cup \{q\}, \Delta \rangle,
\quad
\Delta = \Delta_1 \cup \{ \langle q, \lambda, i \rangle \mid i \in I_1 \} \cup \{ \langle f, \lambda, q \rangle \mid f \in F_1 \}.
\]
No old transition enters $q$, so only the empty word ends there without
reading a word of $L(\mathcal{A}_1)$.

The three constructions in Rascal (\texttt{rename} prefixes the states of each
automaton so that they are disjoint):
\coderange{Rascal}{rascal/src/Automata.rsc}{combine}

\paragraph{Product construction.} Intersection and difference cannot be built
with $\lambda$-transitions, because the automaton has to read the word with both
automata at the same time. For two DFAs $\mathcal{A}_1 = \langle Q_1, \Sigma,
q_{I_1}, F_1, \delta_1 \rangle$ and $\mathcal{A}_2 = \langle Q_2, \Sigma, q_{I_2},
F_2, \delta_2 \rangle$, the product automaton has pairs of states:
\[
\langle Q_1 \times Q_2, \Sigma, (q_{I_1}, q_{I_2}), F, \delta \rangle,
\qquad
\delta((p, q), \sigma) = (\delta_1(p, \sigma), \delta_2(q, \sigma)).
\]
Only the final states change with the operation:
\begin{center}
\begin{tabular}{l l}
  \toprule
  Operation & $F$ \\
  \midrule
  $L_1 \cup L_2$          & $\{ (p, q) \mid p \in F_1 \text{ or } q \in F_2 \}$ \\
  $L_1 \cap L_2$          & $\{ (p, q) \mid p \in F_1 \text{ and } q \in F_2 \}$ \\
  $L_1 \setminus L_2$     & $\{ (p, q) \mid p \in F_1 \text{ and } q \notin F_2 \}$ \\
  \bottomrule
\end{tabular}
\end{center}
Both automata must be complete DFAs: if one of them got stuck, the pair would
have no next state. That is why the set construction keeps the empty set as a
trap state.

Applying the set construction to $L_b$ gives two states: $N = \{W\}$ (no $b$
yet) and $Y = \{W, Z\}$ (a $b$ was read). Applying it to $L_{2a}$ gives three:
$s_0 = \{A\}$, $s_1 = \{A, B\}$ (the word ends in one $a$) and $s_2 = \{A, B,
C\}$ (it ends in $aa$). Figure~\ref{fig:product} shows their product. For
$L_b \cap L_{2a}$ the only final state is $(Y, s_2)$; for $L_{2a} \setminus L_b$
it is $(N, s_2)$; for the union it is every pair with $Y$ or $s_2$.

\begin{figure}[h]
\centering
\begin{tikzpicture}[node distance=2.6cm, auto, every state/.style={draw, minimum size=11mm, inner sep=1pt, font=\small}]
  \node[state, initial] (n0) {$N,s_0$};
  \node[state, right=of n0] (n1) {$N,s_1$};
  \node[state, right=of n1] (n2) {$N,s_2$};
  \node[state, below=2cm of n0] (y0) {$Y,s_0$};
  \node[state, right=of y0] (y1) {$Y,s_1$};
  \node[state, accepting, right=of y1] (y2) {$Y,s_2$};
  \path[->]
    (n0) edge node {a} (n1)
    (n1) edge node {a} (n2)
    (n2) edge[loop right] node {a} ()
    (y0) edge[bend left=15] node {a} (y1)
    (y1) edge node {a} (y2)
    (y2) edge[loop right] node {a} ()
    (y0) edge[loop left] node {b} ()
    (n0) edge node[swap] {b} (y0)
    (n1) edge node[pos=0.25] {b} (y0)
    (n2) edge node[pos=0.15, above] {b} (y0)
    (y1) edge[bend left=15] node {b} (y0)
    (y2) edge[bend left=35] node {b} (y0);
\end{tikzpicture}
\caption{Product of the DFAs for $L_b$ and $L_{2a}$. The final state shown is the one for $L_b \cap L_{2a}$.}
\label{fig:product}
\end{figure}

In Rascal the product is built by searching the reachable pairs, and a
function decides which pairs are final:
\coderange{Rascal}{rascal/src/Automata.rsc}{product}
The tests compare each combination of $L_b$ and $L_{2a}$ with its definition
on every word of length up to 8. They also confirm that the star from the
slide accepts $a$ and the corrected one does not.

\subsection{Exercises}

\begin{exercise}
Define an automaton that accepts the words over $\{0, 1\}$ with an even number
of 0s and an even number of 1s.
\end{exercise}

\paragraph{Solution.} The automaton has to remember two parities, so it needs
four states: EE, OE, EO and OO, where the first letter is the parity of the
0s and the second the parity of the 1s (Figure~\ref{fig:eveneven}). A 0
flips the first letter and a 1 flips the second. EE is initial and final.

\begin{figure}[h]
\centering
\begin{tikzpicture}[node distance=2.8cm, auto]
  \node[state, initial, accepting] (EE) {EE};
  \node[state, right=of EE] (OE) {OE};
  \node[state, below=1.8cm of EE] (EO) {EO};
  \node[state, right=of EO] (OO) {OO};
  \path[->]
    (EE) edge[bend left=15] node {0} (OE)
    (OE) edge[bend left=15] node {0} (EE)
    (EO) edge[bend left=15] node {0} (OO)
    (OO) edge[bend left=15] node {0} (EO)
    (EE) edge[bend left=15] node {1} (EO)
    (EO) edge[bend left=15] node {1} (EE)
    (OE) edge[bend left=15] node {1} (OO)
    (OO) edge[bend left=15] node {1} (OE);
\end{tikzpicture}
\caption{Even number of 0s and even number of 1s.}
\label{fig:eveneven}
\end{figure}

\begin{exercise}
Define an automaton that accepts the words of the language of lists of
Section~\ref{sec:grammars}: \texttt{[]}, \texttt{1|[]}, \texttt{1|2|3|[]}, but not
\texttt{[1,2,3]} or \texttt{1|[]|2|[]}.
\end{exercise}

\paragraph{Solution.} The lecture's automaton is in Figure~\ref{fig:lists-dfa}.
From $q_0$ the word either starts the empty list with \texttt{[} or starts a
number. In $q_3$ the automaton reads the rest of the number; \texttt{|} moves it
to $q_4$, where another number or the final \texttt{[]} may follow. Only
$q_2$ is final, and $q_2$ has no outgoing transitions, so nothing can follow
\texttt{[]}. The second version on the slide adds a state \textit{err} and
sends every missing transition there, which turns the automaton into a
complete DFA.

\begin{figure}[h]
\centering
\begin{tikzpicture}[node distance=2.6cm, auto]
  \node[state, initial] (q0) {$q_0$};
  \node[state, right=of q0] (q1) {$q_1$};
  \node[state, accepting, right=of q1] (q2) {$q_2$};
  \node[state, below=1.8cm of q0] (q3) {$q_3$};
  \node[state, right=of q3] (q4) {$q_4$};
  \path[->]
    (q0) edge node {[} (q1)
    (q1) edge node {]} (q2)
    (q0) edge node[swap] {0..9} (q3)
    (q3) edge[loop left] node {0..9} ()
    (q3) edge[bend left=15] node {|} (q4)
    (q4) edge[bend left=15] node {0..9} (q3)
    (q4) edge node[swap] {[} (q1);
\end{tikzpicture}
\caption{Automaton for the language of lists.}
\label{fig:lists-dfa}
\end{figure}

The Rascal version builds the transitions on digits with a comprehension:
\coderange{Rascal}{rascal/src/Automata.rsc}{lists}

\begin{exercise}
Write an NFA over $\{a, b\}$ that recognizes the words containing the pattern
$bb$ or $bab$ exactly once (occurrences counted from left to right, without
overlapping).
\end{exercise}

\paragraph{Solution.} The lecture's automaton has initial state E, which loops on
$a$ until a $b$ arrives. From A, a second $b$ completes $bb$ (through B and a
$\lambda$ to C) and $a$ then $b$ completes $bab$ (through D). From C on, the
automaton accepts while no second pattern appears: C loops on $a$, a single
$b$ goes to F, $bb$ ends in the trap state H, and $ba$ goes to G, where $b$
would complete a second $bab$. The finals are C, F and G.

The automaton on the slide is missing one transition. After $b$ and $a$ it is
in D, and D only has a transition on $b$. On the word $baabb$ the run is
E, A, D and then stuck, so the word is rejected even though it contains
$bb$ exactly once. The fix is $\langle D, a, E \rangle$: after $baa$ no pattern has
started, so the automaton goes back to the beginning (Figure~\ref{fig:bbbab}).

\begin{figure}[h]
\centering
\begin{tikzpicture}[node distance=2.2cm, auto]
  \node[state, initial] (E) {E};
  \node[state, right=of E] (A) {A};
  \node[state, right=of A] (B) {B};
  \node[state, below=1.5cm of A] (D) {D};
  \node[state, accepting, below=1.5cm of B] (C) {C};
  \node[state, accepting, right=of C] (F) {F};
  \node[state, right=of F] (H) {H};
  \node[state, accepting, below=1.5cm of F] (G) {G};
  \path[->]
    (E) edge[loop above] node {a} ()
    (E) edge node {b} (A)
    (A) edge node {b} (B)
    (A) edge node[swap] {a} (D)
    (B) edge node {$\lambda$} (C)
    (D) edge node[swap] {b} (C)
    (D) edge[dashed] node {a} (E)
    (C) edge[loop below] node {a} ()
    (C) edge node {b} (F)
    (F) edge node {b} (H)
    (F) edge node {a} (G)
    (G) edge node[swap] {b} (H)
    (G) edge[bend left=20] node {a} (C);
\end{tikzpicture}
\caption{Exactly one $bb$ or $bab$. The dashed transition is missing on the slide.}
\label{fig:bbbab}
\end{figure}

\begin{exercise}
Write an NFA that recognizes $(ab \mid aab \mid aba)^*$.
\end{exercise}

\paragraph{Solution.} With transitions over words, the first step is a state E
that is initial and final, with three loops labeled $ab$, $aab$ and $aba$.
Splitting the words into single symbols and sharing the common prefix $a$
gives Figure~\ref{fig:abstar}: from E an $a$ leads to X; from X, $b$ returns
to E ($ab$), $a$ then $b$ returns to E ($aab$), and $b$ then $a$ returns to E
($aba$). The automaton is nondeterministic only in X, which has two
transitions on $b$.

The automaton on the slide puts a loop on $a$ and $b$ in E, makes C the only
final state and adds $\lambda$-transitions from A and D to C. With that loop the
automaton accepts every word that ends in $b$ or $ba$, for example $b$ and
$bbb$, which are not in the language, and it rejects the empty word, which is.

\begin{figure}[h]
\centering
\begin{tikzpicture}[node distance=2.3cm, auto]
  \node[state, initial, accepting] (E) {E};
  \node[state, right=of E] (X) {X};
  \node[state, above right=0.9cm and 1.8cm of X] (Z) {Z};
  \node[state, below right=0.9cm and 1.8cm of X] (Y) {Y};
  \path[->]
    (E) edge[bend left=15] node {a} (X)
    (X) edge[bend left=15] node {b} (E)
    (X) edge node {a} (Z)
    (X) edge node[swap] {b} (Y)
    (Z) edge[bend right=55] node[swap] {b} (E)
    (Y) edge[bend left=55] node {a} (E);
\end{tikzpicture}
\caption{NFA for $(ab \mid aab \mid aba)^*$.}
\label{fig:abstar}
\end{figure}

\coderange{Rascal}{rascal/src/Automata.rsc}{abstar}

\begin{exercise}
Compute the $\lambda$-closure of each state of the automaton with transitions
$A \xrightarrow{\lambda} B$, $A \xrightarrow{\lambda} C$, $B \xrightarrow{a} B$,
$C \xrightarrow{a} D$, $D \xrightarrow{b} E$, $E \xrightarrow{\lambda} C$, with
initial state A and final states B and E.
\end{exercise}

\paragraph{Solution.} $\mathit{cl}_\lambda(A) = \{A, B, C\}$,
$\mathit{cl}_\lambda(B) = \{B\}$, $\mathit{cl}_\lambda(C) = \{C\}$,
$\mathit{cl}_\lambda(D) = \{D\}$ and $\mathit{cl}_\lambda(E) = \{E, C\}$. Eliminating
the $\lambda$-transitions, A becomes final (its closure contains B), and every
transition that enters E also enters C: $D \xrightarrow{b} C$ is added.

\begin{exercise}
\label{ex:neither}
Build a DFA over $\{c, a, t, h\}$ for the words that contain neither
\texttt{cat} nor \texttt{hat}.
\end{exercise}

\paragraph{Solution.} This is the intersection of two languages, so the
product construction applies, but a direct construction is shorter. Both
words end in \texttt{at}, so the automaton only needs to remember whether the
last symbol was \texttt{c} or \texttt{h} (state 1) or whether the last two
symbols were \texttt{ca} or \texttt{ha} (state 2). From state 0, \texttt{c} and
\texttt{h} go to 1 and \texttt{a}, \texttt{t} stay in 0. From 1, \texttt{a} goes
to 2, \texttt{c} and \texttt{h} stay in 1 and \texttt{t} goes back to 0. From 2,
\texttt{t} goes to a trap state, \texttt{c} and \texttt{h} go to 1 and \texttt{a}
goes back to 0. States 0, 1 and 2 are final.
